\documentclass[11pt]{article}
\usepackage{amsmath,amssymb,amsthm}
\usepackage{mathpartir}
\usepackage[margin=2.5cm]{geometry}

\newcommand{\Ctx}{\mathsf{Ctx}}
\newcommand{\wf}[1]{#1\;\mathsf{ctx}}
\newcommand{\hastype}[3]{#1 \vdash #2 : #3}
\newcommand{\conveq}[4]{#1 \vdash #2 = #3 : #4}
\newcommand{\Pitype}[2]{\Pi(#1)#2}
\newcommand{\Lamterm}[2]{\lambda(#1)#2}
\newcommand{\App}[2]{#1\;#2}
\newcommand{\substone}[2]{#1[#2]}
\newcommand{\wkn}[1]{#1^+}
\newcommand{\U}{\mathsf{U}}

\newtheorem{theorem}{Theorem}
\newtheorem{corollary}[theorem]{Corollary}

\title{Typing Rules and Pi Injectivity\\
  for Dependent Type Theory with $\U:\U$}
\author{}
\date{}

\begin{document}
\maketitle

\section{Judgments}

We work with three mutually defined judgment forms:
\begin{align*}
  &\wf{\Gamma}                  &&\text{$\Gamma$ is a well-formed context}\\
  &\hastype{\Gamma}{M}{A}       &&\text{$M$ has type $A$ in context $\Gamma$}\\
  &\conveq{\Gamma}{M}{N}{A}     &&\text{$M$ and $N$ are definitionally equal at type $A$}
\end{align*}
Contexts are snoc-lists: $\Gamma ::= \cdot \mid \Gamma, x{:}A$.

\section{Well-formed contexts}

\begin{mathpar}
  \inferrule{ }{\wf{\cdot}}
  \and
  \inferrule{\hastype{\Gamma}{A}{\U}}{\wf{\Gamma, x{:}A}}
\end{mathpar}

\section{Typing rules}

\begin{mathpar}
  \inferrule[\textsc{Var}]
    {\wf{\Gamma} \\ (x{:}A) \in \Gamma}
    {\hastype{\Gamma}{x}{A}}
  \and
  \inferrule[\textsc{Conv}]
    {\hastype{\Gamma}{M}{A} \\
     \conveq{\Gamma}{A}{B}{\U} \\
     \hastype{\Gamma}{B}{\U}}
    {\hastype{\Gamma}{M}{B}}
  \and
  \inferrule[\textsc{U}]
    {\wf{\Gamma}}
    {\hastype{\Gamma}{\U}{\U}}
  \and
  \inferrule[\textsc{Pi}]
    {\hastype{\Gamma}{A}{\U} \\
     \hastype{\Gamma, x{:}A}{B}{\U}}
    {\hastype{\Gamma}{\Pitype{x{:}A}{B}}{\U}}
  \and
  \inferrule[\textsc{Lam}]
    {\hastype{\Gamma}{A}{\U} \\
     \hastype{\Gamma, x{:}A}{B}{\U} \\
     \hastype{\Gamma, x{:}A}{M}{B}}
    {\hastype{\Gamma}{\Lamterm{x{:}A}{M}}{\Pitype{x{:}A}{B}}}
  \and
  \inferrule[\textsc{App}]
    {\hastype{\Gamma}{A}{\U} \\
     \hastype{\Gamma, x{:}A}{B}{\U} \\
     \hastype{\Gamma}{f}{\Pitype{x{:}A}{B}} \\
     \hastype{\Gamma}{a}{A}}
    {\hastype{\Gamma}{\App{f}{a}}{\substone{B}{a}}}
\end{mathpar}

\section{Conversion rules}

\begin{mathpar}
  \inferrule[\textsc{Refl}]
    {\hastype{\Gamma}{M}{A}}
    {\conveq{\Gamma}{M}{M}{A}}
  \and
  \inferrule[\textsc{Sym}]
    {\conveq{\Gamma}{M}{N}{A}}
    {\conveq{\Gamma}{N}{M}{A}}
  \and
  \inferrule[\textsc{Trans}]
    {\conveq{\Gamma}{M}{N}{A} \\
     \conveq{\Gamma}{N}{P}{A}}
    {\conveq{\Gamma}{M}{P}{A}}
  \and
  \inferrule[\textsc{Conv-Eq}]
    {\conveq{\Gamma}{M}{N}{A} \\
     \conveq{\Gamma}{A}{B}{\U} \\
     \hastype{\Gamma}{B}{\U}}
    {\conveq{\Gamma}{M}{N}{B}}
  \and
  \inferrule[\textsc{Beta}]
    {\hastype{\Gamma}{A}{\U} \\
     \hastype{\Gamma, x{:}A}{B}{\U} \\
     \hastype{\Gamma, x{:}A}{M}{B} \\
     \hastype{\Gamma}{a}{A}}
    {\conveq{\Gamma}
       {\App{(\Lamterm{x{:}A}{M})}{a}}
       {\substone{M}{a}}
       {\substone{B}{a}}}
  \and
  \inferrule[\textsc{Pi-Cong}]
    {\conveq{\Gamma}{A}{A'}{\U} \\
     \conveq{\Gamma, x{:}A}{B}{B'}{\U}}
    {\conveq{\Gamma}
       {\Pitype{x{:}A}{B}}
       {\Pitype{x{:}A'}{B'}}
       {\U}}
  \and
  \inferrule[\textsc{Funext}]
    {\hastype{\Gamma}{A}{\U} \\
     \conveq{\Gamma, x{:}A}{\App{\wkn{f}}{x}}{\App{\wkn{g}}{x}}{B} \\
     \hastype{\Gamma}{f}{\Pitype{x{:}A}{B}} \\
     \hastype{\Gamma}{g}{\Pitype{x{:}A}{B}}}
    {\conveq{\Gamma}{f}{g}{\Pitype{x{:}A}{B}}}
  \and
  \inferrule[\textsc{App-Fun}]
    {\hastype{\Gamma}{A}{\U} \\
     \hastype{\Gamma, x{:}A}{B}{\U} \\
     \conveq{\Gamma}{f}{f'}{\Pitype{x{:}A}{B}} \\
     \hastype{\Gamma}{a}{A}}
    {\conveq{\Gamma}
       {\App{f}{a}}
       {\App{f'}{a}}
       {\substone{B}{a}}}
  \and
  \inferrule[\textsc{App-Arg}]
    {\hastype{\Gamma}{A}{\U} \\
     \hastype{\Gamma, x{:}A}{B}{\U} \\
     \hastype{\Gamma}{f}{\Pitype{x{:}A}{B}} \\
     \conveq{\Gamma}{a}{a'}{A}}
    {\conveq{\Gamma}
       {\App{f}{a}}
       {\App{f}{a'}}
       {\substone{B}{a}}}
\end{mathpar}

\section{Head reduction}

We define head reduction $M \leadsto N$ as the compatible closure
of $\beta$-reduction at the head:
\begin{mathpar}
  \inferrule[\textsc{HR-Beta}]
    { }
    {\App{(\Lamterm{x{:}A}{M})}{a} \leadsto \substone{M}{a}}
  \and
  \inferrule[\textsc{HR-App}]
    {M \leadsto M'}
    {\App{M}{a} \leadsto \App{M'}{a}}
\end{mathpar}

\section{Main result: Pi injectivity (head reduction)}

\begin{theorem}[Corollary 6, part 1; Coquand--Huber 2018, p.\,661]
  If\/ $\conveq{\Gamma}{A_0}{\Pitype{x{:}B_1}{F_1}}{\U}$,
  then there exist $B_0$, $F_0$ such that
  \[
    A_0 \leadsto^* \Pitype{x{:}B_0}{F_0}.
  \]
\end{theorem}

\begin{proof}[Proof outline]
  The proof proceeds by domain-theoretic adequacy.
  \begin{enumerate}
    \item From $\conveq{\Gamma}{A_0}{\Pi B_1\,F_1}{\U}$, extract
      $\hastype{\Gamma}{A_0}{\U}$ by presupposition
      (\textsc{SubstitutionLemma}).
    \item Apply the soundness theorem (\textsc{convSound'}) at the
      \emph{bottom environment} $\rho_\bot$ (mapping every variable to~$\bot$)
      to obtain a bidirectional evaluation transfer: for every finite
      approximation~$u$,
      \[
        \mathsf{EvalRel}(\Pi B_1\,F_1, \rho_\bot, u)
        \;\Longrightarrow\;
        \mathsf{EvalRel}(A_0, \rho_\bot, u).
      \]
    \item Instantiate with $u = \mathsf{PiCode}(\bot, \mathsf{nil})$,
      which is coherent because
      $\mathsf{CoherentFunTail}(\mathsf{nil}) = \top$.
      The evaluation $\mathsf{EvalRel}(\Pi B_1\,F_1, \rho_\bot,
      \mathsf{PiCode}(\bot,\mathsf{nil}))$ is constructible trivially
      (domain~$\bot$, empty function graph, vacuous selection body).
      Transfer gives $\mathsf{EvalRel}(A_0, \rho_\bot,
      \mathsf{PiCode}(\bot,\mathsf{nil}))$.
    \item Apply the bundled adequacy theorem (\textsc{adequacySub2})
      with the identity substitution at~$\rho_\bot$ to obtain
      \[
        \mathsf{Val2}\;\Gamma\;A_0\;\U\;
        (\mathsf{PiCode}(\bot,\mathsf{nil}))\;\mathsf{UCode}.
      \]
    \item By definition, $\mathsf{Val2}$ at $(\mathsf{PiCode}(\bot,
      \mathsf{nil}), \mathsf{UCode})$ unfolds to
      $\mathsf{ValTyPi2}\;\Gamma\;A_0\;\bot\;\mathsf{nil}$, which is a
      $\Sigma$-type whose third component is a contextual reduction
      $\mathsf{Red}\;\Gamma\;A_0\;(\Pi\,B_0\,F_0)\;\U$.
      Extract the underlying head reduction
      $A_0 \leadsto^* \Pi\,B_0\,F_0$.
  \end{enumerate}
  The entire proof is formalised in Agda with \texttt{--without-K
  --exact-split}, no \texttt{--type-in-type}, and no postulates
  (in the files on which \texttt{PiInjectivity.agda} depends).
\end{proof}

\begin{thebibliography}{1}
\bibitem{CH2018}
  T.~Coquand and S.~Huber,
  ``An adequacy theorem for dependent type theory,''
  \emph{Archive for Mathematical Logic}, vol.~57, pp.~649--676, 2018.
\end{thebibliography}

\end{document}
